% !PRINTED 62
\weberpage{62}

Die Function $\Phi(x,\xi)$ wird die \gesperrt{Polare} von $\Phi$ genannt. Sie ist linear und homogen sowohl in Beziehung auf die $x$, wie in Beziehung auf die $\xi$.

Sie kann ausgedrückt werden durch
\begin{equation}
 \Phi(x,\xi)=2\sum a_{ik}\xi_i x_k
\tag{14}
\end{equation}
und genügt der Bedingung
\begin{equation}
 \Phi(x,x)=2\Phi(x).
\tag{15}
\end{equation}

\begin{center}
{\large \S.\ 17.}

\vspace{0.5em}
{\large\gesperrt{Das Euler'sche Theorem über homogene Functionen.}}
\end{center}

Aus den vorstehenden Entwickelungen lässt sich mit Leichtigkeit ein Fundamentalsatz über homogene Functionen herleiten, der von Euler entdeckt und nach ihm benannt ist.

Wir erhalten ihn am einfachsten aus der Formel (4) des vorigen Paragraphen, wenn wir mit $t$ eine beliebige Veränderliche bezeichnen,
\begin{equation}
 \xi_1=tx_1,\quad \xi_2=tx_2,\ldots,\quad \xi_m=tx_m
\tag{1}
\end{equation}
setzen und dann die Fundamentalformel \S~15 (1) für die homogenen Functionen anwenden. Wir erhalten so zunächst:
\begin{equation}
 (1+t)^n\Phi(x)
 =
 \sum { (tx_1)^{\alpha_1}(tx_2)^{\alpha_2}\cdots(tx_m)^{\alpha_m}
 \over \Pi(\alpha_1)\Pi(\alpha_2)\cdots\Pi(\alpha_m)}
 D_{\alpha_1,\alpha_2,\ldots,\alpha_m}\Phi.
\tag{2}
\end{equation}
Wendet man auf der linken Seite von (2) den binomischen Satz an, und setzt dann die Coëfficienten gleich hoher Potenzen von $t$ beiderseits einander gleich, so ergiebt sich für jedes $\nu=1,2,\ldots,n$:
\begin{equation}
 {\Pi(n)\over\Pi(n-\nu)}\Phi(x_1,x_2,\ldots,x_m)
 =
 \sum^{\alpha}
 {\Pi(\nu)\over\Pi(\alpha_1)\Pi(\alpha_2)\cdots\Pi(\alpha_m)}
 x_1^{\alpha_1}x_2^{\alpha_2}\cdots x_m^{\alpha_m}
 D_{\alpha_1,\alpha_2,\ldots,\alpha_m}\Phi,
\tag{3}
\end{equation}
worin sich die Summe auf alle der Bedingung
\begin{equation}
 \alpha_1+\alpha_2+\cdots+\alpha_m=\nu
\tag{4}
\end{equation}
genügenden Werthsysteme der $\alpha$ erstreckt.

In dieser Form ist das zu erweisende Theorem in seiner Allgemeinheit enthalten. Für den besonderen Fall $\nu=1$ erhalten wir die Formel
%CONT>>
